% ============================================================
\chapter{D\'eploiement de Mod\`eles --- REST APIs, FastAPI, Streamlit}
\label{ch:deploiement}
\index{d\'eploiement}
% ============================================================

\section{Patterns de d\'eploiement}
\label{sec:deployment-patterns}
\index{d\'eploiement!patterns}

Un mod\`ele entra\^in\'e n'a de valeur que s'il est accessible aux
utilisateurs ou aux syst\`emes aval. Trois grands patterns de d\'eploiement
existent~:

\begin{definition}[Inf\'erence batch]
L'\textbf{inf\'erence batch}\index{inf\'erence batch} consiste \`a ex\'ecuter
les pr\'edictions sur un lot de donn\'ees \`a intervalles r\'eguliers
(par exemple, chaque nuit). Les r\'esultats sont stock\'es dans une base
de donn\'ees ou un fichier.
\end{definition}

\begin{definition}[Inf\'erence en ligne]
L'\textbf{inf\'erence en ligne}\index{inf\'erence en ligne} (ou temps r\'eel)
r\'epond \`a des requ\^etes individuelles via une API REST ou gRPC, avec
une latence typique de quelques millisecondes \`a quelques secondes.
\end{definition}

\begin{definition}[Inf\'erence en streaming]
L'\textbf{inf\'erence en streaming}\index{inf\'erence streaming} traite un
flux continu de donn\'ees (Kafka, Flink) et produit des pr\'edictions en
continu, adapt\'ee aux cas o\`u les donn\'ees arrivent en temps r\'eel.
\end{definition}

\begin{center}
\begin{tabular}{@{}lccc@{}}
\toprule
\textbf{Crit\`ere} & \textbf{Batch} & \textbf{En ligne} & \textbf{Streaming} \\
\midrule
Latence & Minutes--heures & Millisecondes & Secondes \\
D\'ebit & Tr\`es \'elev\'e & Moyen & \'Elev\'e \\
Complexit\'e & Faible & Moyenne & \'Elev\'ee \\
Cas d'usage & Scoring nocturne & API web/mobile & D\'etection fraude \\
Infra typique & Cron + Spark & FastAPI + K8s & Kafka + Flink \\
\bottomrule
\end{tabular}
\end{center}

\begin{remarque}
La majorit\'e des projets ML en entreprise commencent par du batch, puis
\'evoluent vers du temps r\'eel quand le besoin m\'etier le justifie. Ne
sur-ing\'eniez pas d\`es le d\'epart.
\end{remarque}

\section{Concepts REST API pour le ML}
\label{sec:rest-api}
\index{REST API}

Une API REST expose le mod\`ele via des endpoints HTTP. Les conventions~:

\begin{formules}[title=\textbf{Endpoints typ\-iques d'une API ML}]
\begin{tabular}{@{}llp{6.5cm}@{}}
\toprule
\textbf{M\'ethode} & \textbf{Endpoint} & \textbf{Description} \\
\midrule
\cli{GET} & \cli{/health} & V\'erification que le service est actif \\
\cli{GET} & \cli{/model/info} & M\'etadonn\'ees du mod\`ele (version,
  m\'etriques) \\
\cli{POST} & \cli{/predict} & Pr\'ediction sur une ou plusieurs
  instances \\
\cli{POST} & \cli{/predict/batch} & Pr\'ediction batch (fichier CSV) \\
\bottomrule
\end{tabular}
\end{formules}

\section{FastAPI~: servir un mod\`ele ML}
\label{sec:fastapi}
\index{FastAPI}

\textbf{FastAPI}\index{FastAPI} est un framework Python moderne, rapide
(bas\'e sur Starlette et Pydantic), avec documentation automatique OpenAPI
et validation des donn\'ees int\'egr\'ee.

\subsection{Sch\'emas Pydantic}
\index{Pydantic}

\begin{codeblock}[title=\textbf{schemas.py --- Sch\'emas de requ\^ete et r\'eponse}]
\begin{minted}{python}
"""Pydantic schemas for request/response validation."""
from pydantic import BaseModel, Field


class PredictionRequest(BaseModel):
    """Schema for a single prediction request."""
    sepal_length: float = Field(..., gt=0, description="Sepal length in cm")
    sepal_width: float = Field(..., gt=0, description="Sepal width in cm")
    petal_length: float = Field(..., gt=0, description="Petal length in cm")
    petal_width: float = Field(..., gt=0, description="Petal width in cm")

    model_config = {"json_schema_extra": {
        "examples": [
            {
                "sepal_length": 5.1,
                "sepal_width": 3.5,
                "petal_length": 1.4,
                "petal_width": 0.2,
            }
        ]
    }}


class PredictionResponse(BaseModel):
    """Schema for prediction response."""
    prediction: int = Field(..., description="Predicted class (0, 1, or 2)")
    class_name: str = Field(..., description="Human-readable class name")
    probabilities: list[float] = Field(
        ..., description="Probability for each class"
    )


class HealthResponse(BaseModel):
    """Schema for health check response."""
    status: str
    model_loaded: bool
    model_version: str
\end{minted}
\end{codeblock}

\subsection{Application FastAPI compl\`ete}

\begin{codeblock}[title=\textbf{app.py --- API de serving ML compl\`ete}]
\begin{minted}{python}
"""FastAPI ML serving application."""
import joblib
import numpy as np
from contextlib import asynccontextmanager
from fastapi import FastAPI, HTTPException
from schemas import PredictionRequest, PredictionResponse, HealthResponse

# Variable globale pour le modele
model = None
MODEL_PATH = "models/rf_model.joblib"
MODEL_VERSION = "1.0.0"
CLASS_NAMES = ["setosa", "versicolor", "virginica"]


@asynccontextmanager
async def lifespan(app: FastAPI):
    """Load model on startup, cleanup on shutdown."""
    global model
    try:
        model = joblib.load(MODEL_PATH)
        print(f"Model loaded from {MODEL_PATH}")
    except FileNotFoundError:
        print(f"WARNING: Model not found at {MODEL_PATH}")
    yield
    model = None
    print("Model unloaded")


app = FastAPI(
    title="ML Prediction API",
    description="API de prediction pour le modele Iris",
    version=MODEL_VERSION,
    lifespan=lifespan,
)


@app.get("/health", response_model=HealthResponse)
async def health_check():
    """Health check endpoint."""
    return HealthResponse(
        status="healthy" if model is not None else "degraded",
        model_loaded=model is not None,
        model_version=MODEL_VERSION,
    )


@app.post("/predict", response_model=PredictionResponse)
async def predict(request: PredictionRequest):
    """Generate prediction for a single instance."""
    if model is None:
        raise HTTPException(status_code=503, detail="Model not loaded")

    features = np.array([[
        request.sepal_length,
        request.sepal_width,
        request.petal_length,
        request.petal_width,
    ]])

    prediction = int(model.predict(features)[0])
    probabilities = model.predict_proba(features)[0].tolist()

    return PredictionResponse(
        prediction=prediction,
        class_name=CLASS_NAMES[prediction],
        probabilities=[round(p, 4) for p in probabilities],
    )
\end{minted}
\end{codeblock}

\begin{codeblock}[title=\textbf{Lancer et tester l'API}]
\begin{minted}{bash}
# Installer FastAPI et uvicorn
pip install fastapi uvicorn

# Lancer le serveur
uvicorn app:app --host 0.0.0.0 --port 8000 --reload

# Tester avec curl
curl -X POST http://localhost:8000/predict \
  -H "Content-Type: application/json" \
  -d '{"sepal_length": 5.1, "sepal_width": 3.5,
       "petal_length": 1.4, "petal_width": 0.2}'

# Documentation automatique
# Swagger UI : http://localhost:8000/docs
# ReDoc :      http://localhost:8000/redoc
\end{minted}
\end{codeblock}

\begin{codeoutput}
\begin{minted}{json}
{
  "prediction": 0,
  "class_name": "setosa",
  "probabilities": [0.97, 0.02, 0.01]
}
\end{minted}
\end{codeoutput}

\section{Architecture de d\'eploiement}
\label{sec:deployment-architecture}

\begin{center}
\begin{tikzpicture}[
  box/.style={rectangle, draw, rounded corners, minimum width=2.5cm,
    minimum height=1.2cm, align=center, font=\small},
  arr/.style={-{Stealth[length=2.5mm]}, thick},
  dashed-arr/.style={-{Stealth[length=2.5mm]}, thick, dashed, gray}
]
  % Client
  \node[box, fill=blue!10] (client) {Client\\(Web/Mobile)};

  % Load balancer
  \node[box, fill=yellow!10, right=2cm of client] (lb) {Load\\Balancer};

  % API instances
  \node[box, fill=green!10, right=2cm of lb, yshift=1cm] (api1) {FastAPI\\Instance 1};
  \node[box, fill=green!10, right=2cm of lb, yshift=-1cm] (api2) {FastAPI\\Instance 2};

  % Model store
  \node[box, fill=orange!10, right=2cm of api1, yshift=-1cm] (store) {Model\\Store};

  % Monitoring
  \node[box, fill=red!10, below=1.5cm of lb] (mon) {Monitoring\\(Prometheus)};

  % Arrows
  \draw[arr] (client) -- (lb);
  \draw[arr] (lb) -- (api1);
  \draw[arr] (lb) -- (api2);
  \draw[dashed-arr] (store) -- (api1);
  \draw[dashed-arr] (store) -- (api2);
  \draw[dashed-arr] (api1) -- (mon);
  \draw[dashed-arr] (api2) -- (mon);
\end{tikzpicture}
\end{center}

\section{Streamlit~: construire un d\'emo UI}
\label{sec:streamlit}
\index{Streamlit}

\textbf{Streamlit}\index{Streamlit} permet de cr\'eer rapidement une
interface web interactive pour d\'emontrer un mod\`ele ML, sans conna\^itre
le d\'eveloppement frontend.

\begin{codeblock}[title=\textbf{streamlit\_app.py --- D\'emo interactive}]
\begin{minted}{python}
"""Streamlit demo for the Iris classification model."""
import streamlit as st
import requests
import json

API_URL = "http://localhost:8000"

st.set_page_config(page_title="Iris Classifier", page_icon="🌸")
st.title("Classification Iris")
st.markdown("Entrez les mesures de la fleur pour obtenir une prediction.")

# Sidebar pour les parametres
st.sidebar.header("Parametres de la fleur")
sepal_length = st.sidebar.slider("Longueur sepale (cm)", 4.0, 8.0, 5.1)
sepal_width = st.sidebar.slider("Largeur sepale (cm)", 2.0, 4.5, 3.5)
petal_length = st.sidebar.slider("Longueur petale (cm)", 1.0, 7.0, 1.4)
petal_width = st.sidebar.slider("Largeur petale (cm)", 0.1, 2.5, 0.2)

# Bouton de prediction
if st.button("Predire"):
    payload = {
        "sepal_length": sepal_length,
        "sepal_width": sepal_width,
        "petal_length": petal_length,
        "petal_width": petal_width,
    }

    try:
        response = requests.post(f"{API_URL}/predict", json=payload)
        response.raise_for_status()
        result = response.json()

        st.success(f"Prediction : **{result['class_name']}**")

        # Afficher les probabilites
        st.subheader("Probabilites par classe")
        class_names = ["setosa", "versicolor", "virginica"]
        for name, prob in zip(class_names, result["probabilities"]):
            st.progress(prob, text=f"{name}: {prob:.1%}")

    except requests.exceptions.ConnectionError:
        st.error("Impossible de se connecter a l'API. "
                 "Verifiez que le serveur tourne.")

# Health check
st.sidebar.markdown("---")
if st.sidebar.button("Verifier l'API"):
    try:
        health = requests.get(f"{API_URL}/health").json()
        st.sidebar.json(health)
    except requests.exceptions.ConnectionError:
        st.sidebar.error("API indisponible")
\end{minted}
\end{codeblock}

\begin{codeblock}[title=\textbf{Lancer le dashboard Streamlit}]
\begin{minted}{bash}
# Installer Streamlit
pip install streamlit requests

# Lancer l'application
streamlit run streamlit_app.py --server.port 8501

# Accessible sur http://localhost:8501
\end{minted}
\end{codeblock}

\section{S\'erialisation de mod\`eles}
\label{sec:serialization}
\index{s\'erialisation}

Le choix du format de s\'erialisation affecte la portabilit\'e, la vitesse
d'inf\'erence et la compatibilit\'e inter-frameworks.

\begin{formules}[title=\textbf{Formats de s\'erialisation}]
\begin{tabular}{@{}lp{3.5cm}p{3.5cm}p{3cm}@{}}
\toprule
\textbf{Format} & \textbf{Avantages} & \textbf{Inconv\'enients} & \textbf{Usage} \\
\midrule
\textbf{joblib/pickle} & Simple, natif Python & Python uniquement,
  vuln\'erable & Prototypage, scikit-learn \\
\textbf{ONNX} & Multi-framework, optimis\'e & Pas tous les op\'erateurs &
  Production cross-platform \\
\textbf{TorchScript} & Optimis\'e PyTorch, C++ & PyTorch uniquement &
  Production PyTorch \\
\textbf{SavedModel} & \'Ecosyst\`eme TF complet & TensorFlow uniquement &
  TF Serving, TFLite \\
\bottomrule
\end{tabular}
\end{formules}

\begin{codeblock}[title=\textbf{Export de mod\`eles dans diff\'erents formats}]
\begin{minted}{python}
"""Model serialization examples."""
import joblib
import torch
import onnx
from skl2onnx import convert_sklearn
from skl2onnx.common.data_types import FloatTensorType

# --- joblib (scikit-learn) ---
joblib.dump(model, "model.joblib")
model = joblib.load("model.joblib")

# --- ONNX (scikit-learn) ---
initial_type = [("features", FloatTensorType([None, 4]))]
onnx_model = convert_sklearn(model, initial_types=initial_type)
onnx.save(onnx_model, "model.onnx")

# Inference ONNX
import onnxruntime as ort
session = ort.InferenceSession("model.onnx")
input_name = session.get_inputs()[0].name
result = session.run(None, {input_name: features.astype("float32")})

# --- TorchScript (PyTorch) ---
scripted_model = torch.jit.script(pytorch_model)
scripted_model.save("model.pt")
loaded = torch.jit.load("model.pt")
output = loaded(torch.tensor(features))
\end{minted}
\end{codeblock}

\begin{attention}[title=\textbf{S\'ecurit\'e de pickle/joblib}]
Les fichiers \file{.pkl} et \file{.joblib} peuvent ex\'ecuter du code
arbitraire au chargement. Ne chargez \textbf{jamais} un mod\`ele
s\'erialis\'e provenant d'une source non fiable. En production, pr\'ef\'erez
ONNX ou TorchScript.
\end{attention}

\section{Comparaison des frameworks de serving}
\label{sec:serving-comparison}
\index{serving}

\begin{center}
\begin{tabular}{@{}lcccc@{}}
\toprule
\textbf{Crit\`ere} & \textbf{FastAPI} & \textbf{Flask} & \textbf{TorchServe} & \textbf{Triton} \\
\midrule
Async natif & Oui & Non & Oui & Oui \\
Validation auto & Pydantic & Manuelle & Manuelle & Manuelle \\
Doc OpenAPI & Automatique & Manuelle & Non & Non \\
Multi-framework & Oui & Oui & PyTorch & Tous \\
GPU optimis\'e & Manuel & Manuel & Oui & Oui \\
Batching & Manuel & Manuel & Automatique & Automatique \\
Monitoring & Manuel & Manuel & Int\'egr\'e & Int\'egr\'e \\
Complexit\'e & Faible & Faible & Moyenne & \'Elev\'ee \\
Cas d'usage & API g\'en\'erale & Legacy & PyTorch prod & Haute perf \\
\bottomrule
\end{tabular}
\end{center}

\begin{remarque}
Pour d\'ebuter, FastAPI est le meilleur compromis entre simplicit\'e et
performance. TorchServe et Triton sont indiqu\'es pour les d\'eploiements
\`a haute performance n\'ecessitant du batching dynamique et une
optimisation GPU avanc\'ee.
\end{remarque}

\section{Dockerfile pour le d\'eploiement}
\label{sec:deploy-dockerfile}

\begin{codeblock}[title=\textbf{Dockerfile pour l'API + Streamlit}]
\begin{minted}{dockerfile}
# === API FastAPI ===
FROM python:3.11-slim AS api

WORKDIR /app
COPY requirements-api.txt .
RUN pip install --no-cache-dir -r requirements-api.txt

COPY src/ ./src/
COPY models/ ./models/
COPY schemas.py app.py ./

RUN useradd --create-home mluser && chown -R mluser:mluser /app
USER mluser

EXPOSE 8000
HEALTHCHECK --interval=30s --timeout=5s --retries=3 \
  CMD curl -f http://localhost:8000/health || exit 1

CMD ["uvicorn", "app:app", "--host", "0.0.0.0", "--port", "8000"]
\end{minted}
\end{codeblock}

\begin{codeblock}[title=\textbf{docker-compose.yml pour API + Streamlit}]
\begin{minted}{yaml}
version: "3.9"

services:
  api:
    build:
      context: .
      dockerfile: Dockerfile
      target: api
    ports:
      - "8000:8000"
    volumes:
      - ./models:/app/models
    healthcheck:
      test: ["CMD", "curl", "-f", "http://localhost:8000/health"]
      interval: 30s
      timeout: 5s
      retries: 3

  streamlit:
    build:
      context: .
      dockerfile: Dockerfile.streamlit
    ports:
      - "8501:8501"
    environment:
      - API_URL=http://api:8000
    depends_on:
      api:
        condition: service_healthy
\end{minted}
\end{codeblock}

\section{Bonnes pratiques de d\'eploiement}
\label{sec:deploy-best-practices}

\begin{bonnepratique}[title=\textbf{R\`egles d'or du d\'eploiement ML}]
\begin{enumerate}
  \item \textbf{Healthcheck obligatoire}~: le endpoint \cli{/health} doit
    v\'erifier que le mod\`ele est charg\'e et fonctionnel.
  \item \textbf{Versionner les mod\`eles}~: inclure la version dans
    l'endpoint \cli{/model/info} et dans les r\'eponses.
  \item \textbf{Valider les entr\'ees}~: utiliser Pydantic pour rejeter
    les requ\^etes malform\'ees avant l'inf\'erence.
  \item \textbf{Logging structur\'e}~: chaque pr\'ediction doit \^etre
    logg\'ee avec les features d'entr\'ee, la pr\'ediction et la latence.
  \item \textbf{Timeout et retry}~: configurer des timeouts c\^ot\'e
    client et serveur.
  \item \textbf{S\'eparer API et UI}~: FastAPI et Streamlit dans des
    conteneurs s\'epar\'es.
  \item \textbf{Load testing}~: utiliser \pkg{locust} ou \pkg{wrk} pour
    mesurer la capacit\'e avant la mise en production.
\end{enumerate}
\end{bonnepratique}

\begin{antipattern}[title=\textbf{Anti-patterns de d\'eploiement}]
\begin{itemize}
  \item Charger le mod\`ele \`a chaque requ\^ete (utiliser le lifespan).
  \item Pas de validation des entr\'ees (crash sur donn\'ees inattendues).
  \item Retourner des erreurs g\'en\'eriques (toujours <<~500 Internal Server
    Error~>>) sans d\'etails utiles.
  \item Pas de healthcheck (le load balancer envoie du trafic \`a un service
    mort).
  \item D\'eployer sans tests d'int\'egration de l'API.
  \item Ignorer la latence~: ne pas mesurer le temps d'inf\'erence.
\end{itemize}
\end{antipattern}

\section{Mini-projet~: d\'eployer le mod\`ele du cours}
\label{sec:mini-project-ch8}

\begin{miniprojet}[title=\textbf{Mini-projet 8~: API + D\'emo Streamlit}]
En reprenant le mod\`ele entra\^in\'e dans les chapitres pr\'ec\'edents~:
\begin{enumerate}
  \item \textbf{S\'erialisation}~: sauvegardez le mod\`ele en \file{.joblib}
    et en ONNX. Comparez les tailles.
  \item \textbf{API FastAPI}~:
    \begin{itemize}
      \item Cr\'eez les sch\'emas Pydantic adapt\'es \`a votre mod\`ele
      \item Impl\'ementez \cli{/health}, \cli{/model/info} et \cli{/predict}
      \item Ajoutez un logging de chaque pr\'ediction (features, r\'esultat,
        latence)
    \end{itemize}
  \item \textbf{D\'emo Streamlit}~: cr\'eez un dashboard interactif qui~:
    \begin{itemize}
      \item Permet de saisir les features via des sliders
      \item Affiche la pr\'ediction et les probabilit\'es
      \item Affiche un historique des derni\`eres pr\'edictions
    \end{itemize}
  \item \textbf{Conteneurisation}~: cr\'eez un \file{docker-compose.yml}
    avec les services \cli{api} et \cli{streamlit}.
  \item \textbf{Tests}~: \'ecrivez des tests d'int\'egration avec
    \pkg{httpx} ou \cli{TestClient} de FastAPI.
  \item \textbf{Load testing}~: mesurez la latence et le d\'ebit avec
    \pkg{locust} (objectif~: $<$100~ms p95 pour une instance).
\end{enumerate}
\end{miniprojet}

\section{Exercices}
\label{sec:exercises-ch8}

\begin{exercice}[$\bigstar$ --- Premi\`ere API FastAPI]
Cr\'eez une API FastAPI minimaliste qui charge un mod\`ele scikit-learn
entra\^in\'e et expose un endpoint \cli{/predict}. Testez-la avec
\cli{curl} et la documentation Swagger automatique
(\cli{http://localhost:8000/docs}).
\end{exercice}

\begin{exercice}[$\bigstar\bigstar$ --- Validation et gestion d'erreurs]
Am\'eliorez l'API de l'exercice pr\'ec\'edent~:
\begin{enumerate}
  \item Ajoutez des sch\'emas Pydantic avec des contraintes de validation
    (bornes min/max, types stricts)
  \item Impl\'ementez un endpoint \cli{/predict/batch} qui accepte une
    liste d'instances
  \item Ajoutez une gestion d'erreurs compl\`ete~: mod\`ele non charg\'e
    (503), donn\'ees invalides (422), erreur interne (500) avec messages
    clairs
  \item Ajoutez un middleware de logging qui enregistre chaque requ\^ete
    (m\'ethode, endpoint, latence, code de retour)
  \item \'Ecrivez des tests unitaires avec \cli{TestClient} de FastAPI
\end{enumerate}
\end{exercice}

\begin{exercice}[$\bigstar\bigstar\bigstar$ --- D\'eploiement complet avec ONNX]
R\'ealisez un d\'eploiement production-ready~:
\begin{enumerate}
  \item Exportez le mod\`ele en ONNX. Comparez la latence d'inf\'erence
    entre joblib et ONNX Runtime
  \item Cr\'eez une API FastAPI avec support ONNX (\pkg{onnxruntime})
  \item Ajoutez un dashboard Streamlit connect\'e \`a l'API
  \item Conteneurisez le tout avec Docker Compose (API + Streamlit)
  \item R\'edigez un load test avec \pkg{locust}~: mesurez la latence
    p50, p95, p99 et le d\'ebit maximal
  \item Configurez un reverse proxy Nginx devant l'API
  \item Documentez l'architecture dans un sch\'ema (TikZ ou draw.io)
\end{enumerate}
\end{exercice}

\section{Aide-m\'emoire}

\begin{cheatsheet}[title=\textbf{R\'esum\'e du chapitre 8}]
\begin{tabular}{@{}p{4cm}p{8.5cm}@{}}
\toprule
\textbf{Concept} & \textbf{Point cl\'e} \\
\midrule
Patterns & Batch (cron) vs En ligne (API) vs Streaming (Kafka) \\
FastAPI & \cli{uvicorn app:app --port 8000}, Pydantic, async \\
Streamlit & \cli{streamlit run app.py}, sliders, \cli{st.button} \\
S\'erialisation & joblib (proto), ONNX (prod), TorchScript (PyTorch) \\
Healthcheck & \cli{GET /health} avec statut mod\`ele \\
Validation & Pydantic \cli{BaseModel} avec \cli{Field} contraintes \\
Serving avanc\'e & TorchServe (PyTorch), Triton (multi-framework) \\
Anti-pattern & Charger le mod\`ele par requ\^ete, pas de healthcheck \\
\bottomrule
\end{tabular}
\end{cheatsheet}
